\chapter{Book IV: the pocket book, fifty-four years without a gap}

\section{An object unlike the others}

Book IV came to the Whitney folded inside a used envelope, typewritten \q{EDWARD HOPPER / 3 WASHINGTON SQUARE N. / NEW YORK CITY, NEW YORK} and metered at Lynbrook in July 1963 (ref 16225). It is a stationer's pocket memorandum book, \q{MEMORANDUM} stamped on brown cloth, seven and a half inches by four and three-quarters, its pages ruled into a date column, a description column and a money column divided into dollars and cents by two printed red rules. Nobody adds a column and nobody heads one. It begins on 15 November 1913 and ends on 23 March 1967, and between those dates there is no gap. It is the only one of the six volumes that the Whitney dated leaf by leaf, and it is, for that reason, the volume the edition's timeline can point into.

It is also in a hand the edition does not name. The first batch describes \q{a fluent sloping copperplate} unlike either hand recorded from the work books, and every passage in all fourteen batches is attributed to nobody. The entries begin eleven years before Josephine Nivison joined the household. Whether the hand of 1913 and the hand of 1967 are the same, and whether either is Edward Hopper's, is a question the transcription puts and does not answer. What it can say is how the book was kept, because the book has a method and the method is written in colour.

\section{Charges black, receipts red, sums pencil}

On leaves 3 to 6 the text is black and the money figures are red. From leaf 7, in May 1914, the colours swap and hold for fifty-three years: the charge in black, the receipt line in red, \q{Rec'd by check} and its figure written across the whole width of the leaf; and, from leaf 66, a pencil sum at the foot of the leaf, and from leaf 114 a labelled total for the year. The edition records the ink where it has been read off the photographs, and where it has, \q{red money is a receipt and pencil money a sum before any word is read}. This is not a detail of the transcription; it is the accounting system, and a black-and-white reproduction of Book IV would not be a reproduction of it.

The first entry is the model for the first decade:

\begin{leafquote}
1913 \quad Nov. 15th \quad Asc. Sunday Magazines / 2 line drawings / “The Iron Stickpin” \quad 20 00 / “Wayland” \quad 20 00 / Nov. 22nd \quad Rec'd by check \quad \$40 00 \hfill (\lf{Book IV}{3}{17106})
\end{leafquote}

\noindent A client, its items one to a line with their prices, and the cheque a week later. Twenty dollars a line drawing for the Associated Sunday Magazines is the baseline of the volume, and the illustration years are a catalogue of what a commercial artist in New York was paid between 1913 and 1925. The clients are named in full: U.S. Printing and Litho for Eclair film posters at ten and twelve dollars a sheet, with \q{Time watching films 2}; the A. W. Shaw Company, half tones at twenty-five a half page; J. Walter Thompson, \q{Enos Salts} for twenty; Everybody's Magazine, four war drawings, \q{1 Uhlans charging 25 / 1 Infantry attack 25 / 1 German gun in action 40 / 1 French in trenches 25}; Webb Publishing of St Paul, twenty a drawing; the Wells Fargo Messenger, \q{1 wagon poster 50}; Adventure, the long lists at fifteen a drawing, fifteen drawings for \$225 with titles like \q{The Man-Breaker}, \q{Snuffy and the Monster}, \q{The Snagdangle Twins' Last Cleanup}; the Country Gentleman; the Morse Dry Dock Dial, its cover nearly every month from January 1918 at thirty-five, forty, forty-five and fifty dollars, the subjects \q{painter}, \q{man at lathe}, \q{welder}, \q{boiler maker}; Scribner's, a steady \$175 for three drawings from 1922 to 1925; and the Hotel Management cover, sixty-five a month, from June 1924. In the middle of the war sits the largest sum of the first forty leaves, and it has no charge line: \q{U. S. Shipping Board / 1 poster / Smash the Hun / rec'd by check 300} (\lf{Book IV}{35}{18332}). The prize was three hundred dollars, and the book, which had no way to write a receipt without a charge, wrote it anyway.

Two other kinds of income sit among the drawings. \q{Mrs L. Whitins / instruction 7 weeks 105}: he taught, at fifteen dollars a week (\textsc{Book IV}, leaf 67). And the New Republic paid \q{plate Night Shadows 100} and \q{Robert Hallowell / advance royalties 100} for the portfolio of 1924 (\lf{Book IV}{71}{17582}). The last illustration money is a Hotel Management cover billed on 3 December 1925 and paid on Christmas Eve (\lf{Book IV}{77}{16585}). The transcription's own header says November; the leaf says December; the difference is one cover, and it is the kind of thing a second reader is for.

\section{The pictures arrive}

The first painting in the book is two lines, out of date order: \q{Nov 16 Bklyn Museum / The Mansarde Roof 100 / Dec 14 rec'd by check 100} (\lf{Book IV}{66}{17008}), in 1923. The first etching sold is on leaf 49, in 1920, to the Print Makers of Los Angeles: \q{The Monhegan Boat. / @ \$18. less 10\% 16 20}, the first use of the cents column. The prizes of 1923 are written as purchases: \q{Chicago Society of Etchers / prize award the Logan / 1 East Side Interior 25}. The Metropolitan's fifteen prints at ten dollars are on leaf 75, in May 1925. The British Museum paid \q{£1 / rec'd by cash £1} with a pencilled \q{5 00} beneath, the only sterling in the book (\lf{Book IV}{80}{16737}).

And then, on leaf 71, 24 November 1924: \q{Frank K. M. Rehn / 5 Water colors at 33 1/3 \% / 1 100 / 2 100 / 3 100 / 4 100 / 5 100 / '' 24 rec'd by check 500}. The watercolours are unnamed and numbered down the column; the hundred is the net; the series runs to fourteen by January 1925. Rehn had sold the Gloucester watercolours off the walls of his back room, and the pocket book records the money without recording the pictures. From the end of 1925 the works are named: \q{3 Water colors / the Locomotive of D \& G. 133 34 / Adobe Houses 133 34 / La Penitente 133 34 / rec'd 400} (\lf{Book IV}{77}{16585}). The illustration stops the same month. The two trades, which the work books never mention together, hand over to one another on one leaf.

\section{Four ways of writing a commission}

The edition's method note for Book IV divides the volume into four periods by how the commission is written, and the division is the best account we have of how a dealer's statement looked to the person who received it. In 1924--27 the rate is written, \q{at 33 1/3 \%}. In 1927--30 no rate is written at all, only nets whose cents carry the third: 166.67 for a watercolour, 16.67 for a print. In 1931--37 the deductions are named cash sums, \q{less \$49.93}, \q{less .50 for mats}, \q{less photo \$16.}, which is the photographer's bill passed through and not the dealer's rate; the Whitney's purchase of Lombard's House is written in red on the receipt, \q{rec'd by checks 750. / commission to Rehn 250.} (\lf{Book IV}{101}{17933}), and the Circle Theatre, \q{\$1,800}, is the first dollar sign and comma in the money column. From 1938 to 1955 one sale is written down the leaf four or five times: the items at list price, each marked \q{o}, \q{w.c} or \q{E}; a rule across the money column and the sum; \q{less commission}, which names either the deduction or the net, the same words for both; \q{less photos}; the net restated; the red cheque. And from 1955 the entry is the gallery's statement, headed \q{Rec'd from Rehn Gallery}, sometimes with the cheque number, \q{check No 6388 date Feb 14 / Empire Trust Co 7798 43} (\lf{Book IV}{140}{18250}), and sometimes with the pictures counted and not priced: \q{10 water colors / 4 drawings / 3 etchings / 3 oils 16,514 65} (\lf{Book IV}{141}{16957}). By then the book is no longer an invoice book but a record of what arrived.

The prices can be read across the decades. Watercolours at 100 net in 1924, 166.67 in 1927; The City 600 in 1927; \q{Light House Hill (Scheffer) 1000}; Blackwell's Island 1666.67 net in 1929; \q{“Church So. Truro” / “Tables for Ladies”} 4200 gross in 1931; \q{“Room in Brooklyn”} 1200 to Boston; \q{From Williamsburg Bridge 2000} in 1937; \q{o Cape Cod Afternoon 3800 / o. Compartment C Car 293 1000} in 1938; \q{“New York Movie” 2500}; \q{“Night Hawks” 3000} and \q{“Dawn in Pennsylvania” 2000} in 1942; \q{“Gas” 1800} in 1943; \q{“Hotel Lobby” 2500}; \q{Rooms for Tourists oil 3000}; \q{1 oil Rooms by the Sea 4000 00} in 1952; \q{Morning in a City 6000 00}; \q{1 oil “Hotel Room” 5500 00}; \q{Sunlight in a Cafeteria 9000 00} in 1958; \q{November Washington Square / Santa Barbara Museum / Mrs Sterling Morton 10,000.00} in 1960; \q{Girlie Show (Charles Stein) 11000 00}; \q{“People in the Sun” 15000 00}; \q{“Excursion into Philosophy” \$14500.00}; \q{“New York Office” 16500.00 / Lawrence Fleischman}; \q{“Sun in an Empty Room” 25000 00}; \q{balance on Intermission 25000 00}. The largest single cheque in the book is 17,737.84, on 11 January 1958, against Rehn's statement of 28 December 1957; the largest of the 1940s was 4357.76 and of the 1920s 3380.17.

Until 1954 the pocket book names almost no buyers: \q{(Phillips)}, \q{(Mrs Coburn)}, \q{(Scheffer)}, in brackets, and museums as destinations. Leaf 133 is \q{the first leaf in Book IV to name a buyer}, and only in a pencil margin: \q{Roy Neuberger} against the Barber Shop, \q{Jo Hirshhorn} against four pictures. Stephen Clark, the collector of Book III, is not in Book IV at all. From 1961 the buyer is in the charge line, because the gallery's statement now carried it. The pocket book is the dealer's side of the record; the work books are the painter's.

\section{The instalments}

In the 1960s a picture is paid for in parts, and the book writes each part as an entry of its own. \q{Excursion into Philosophy / on acct. 5000 00} among the items of 24 January 1962; then, on 23 June, the picture at \$14,500 with \q{Less payment on acount / '' 1/24/1962 / “Excursion into Philosophy” 3333.34} taken off the net, which is the five thousand less a third. \q{A Woman in the Sun 10000.00} in June 1962; \q{“A Woman in the Sun” 15000.00 / Mr \& Mrs Albert Hackett} in January 1963 with \q{Payment on Account / Hackett --- 6666.67} deducted. \q{Road and Trees 15000 00 / Rehn Gallery / John Clancy purchase} twice on one leaf, 7500 received on account in July 1963 and 2500 as the balance in December, no commission line, the two receipts making two thirds of the price: the dealer bought it for himself (\lf{Book IV}{154}{17869}). \q{Intermission 6,000 00} on account in December 1965; \q{balance on Intermission 25000 00 / less 33 1/3 8333 33 / 16666 67 / Dec. 16, 1965 / Payment on acct 6000 00 / check --- 10666 67} in July 1966 (\lf{Book IV}{158}{16958}). Chair Car: \q{payment on account / “Chair Car” 10,000 00} in December 1965, \q{Balance on Chair Car / recd by check 8,000 00 / to be deposited by John / in Bank of New York on / Jan 1, 1967} in December 1966, and the full price never written anywhere in the book. Read row by row, every one of those pictures was charged twice and the years 1962 and 1963 read 25,000 high; the edition's reader collapses them and says so. A historian using this book for income has to know that the book was not kept for that purpose. It was kept to know what had come in.

\section{What the book says it counts}

Once, in fifty-four years, the book states a rule about itself. Against a Kennedy sale of a Cat Boat print on 14 January 1944, in a smaller upright pencil, \q{Take note} and \q{not included in income.}, with a pencil cross in the money column where a figure would be (\lf{Book IV}{116}{18385}). The transcription calls it \q{the only place in ten batches where the book says what it counts}. Refunds and insurance are left out of the sums, prizes inconsistently; a loan is counted, \q{Insurance on Capt Ed / Staples from Cleveland Mus. 1000} in 1929, the picture that burnt on the train. Money going the other way is written too: \q{Cape \& Vineyard Electric Co / for electric line at Truro 914 00}, in pencil, unreceipted (\lf{Book IV}{143}{18525}), which is the electricity the black notebook's chronicle says they \q{began efforts for} in 1951 and \q{got} in September 1954; \q{Paid \$5000.00 on Fed Income Tax / Feb. 15, 1957} in pencil under the cheque of 16,514.65; \q{Frame for Dauphinee House -48 00}; \q{Relining of “East River” / (cleaning and varnishing) 190 00 / Frame by Kulicke for / “East River” 144 00}.

The year totals she wrote from leaf 114 onward, in pencil, are the check on the whole reading. \q{1942 --- 5857.44}; \q{1949 total 3661.14}; \q{Total 1955} 13406.49 with its last digit cut by the binding; \q{Total 39842.56} for 1962. Where the leaf's own arithmetic closes, the edition's receipts stream reconciles with it to the cent: 1946, 1947, 1949, 1951, 1961, 1962. Where it does not, the difference is named: ten dollars short in 1952, \q{Thirty-six dollars are unaccounted for} in 1958, the thousand-dollar award left out of 1960. The book also moves money between years and says so: seven thousand dollars of 1954 receipted again in 1955 with \q{delayed receipt Re Zucker --- p income Tax} beside it; a cheque of 27 December 1955 \q{Jan 3 deposited}; the Chair Car balance dated 15 December 1966 and entered again on 1 January 1967. These are not errors. They are a cash book keeping cash-basis accounts against a tax year, and the tax year was the reason the totals began to be labelled at all.

\section{The last leaf}

Leaf 158 carries the only writing in the volume the edition cannot place. \q{Dec. 15, 1966 / Balance on Chair Car} is \q{written in blue and then written again in black over the blue, word for word and a little out of register}, and the deposit note \q{in red and then written again in black over the red}. Nothing is added or altered; somebody went over the last entries a second time, and nothing on the leaf says who.

Leaf 159, 1967, holds three entries, each with a bare red figure under the black charge: the 8000 deposited by Clancy on 1 January; the Museum of Modern Art's 100 \q{for reproduction of House by / a Railroad, Illustration / on jacket / The “American Tradition / in Literature”}; and the last entry in Book IV: \q{Mar. 23 Rec'd from Artex Prints / Royalties on Light House a 2 lights / by Artex Junior / recd by check Mar 14. 28 83}, with a red \q{Mar 16 28 23} beneath, \q{receipted twice at two figures a fortnight apart}. Twenty-eight dollars in royalties on a reproduction of a lighthouse painted in 1929, seven weeks before he died. \q{The rest of the leaf is empty.} The back cover: \q{No hand has written on it.}

Nothing in the volume is dated after 23 March 1967. Ninety-two leaves earlier, the same hand had written \q{2 line drawings / “The Iron Stickpin” 20 00}. The whole of a working life is between them, and the book never once says what any of it was for.
